\subsection{生成模的双交换子}

定理 2. —— 生成模是平衡的。

设 A 是一个环，M 是生成 A-模；按定义，存在整数 \(n \geqslant 0\)、M 的元素 \(x_{1}, \ldots, x_{n}\) 及 \(x_{1}^{*}, \ldots, x_{n}^{*}\)（M 的对偶 \(M^{*}\) 的元素），满足 \(\sum_{i=1}^{n} \langle x_{i}, x_{i}^{*} \rangle = 1\)。回忆（II, §4, No.~2, 第 271 页），我们用下述公式定义群同态 \(\theta : M^{*} \otimes_{A} M \to \operatorname{End}_{A}(M)\)：\(\theta(x^{*} \otimes y)(x) = \langle x, x^{*} \rangle y\)。若 u 是 M 的双交换子的一个元素，则它与 \(\operatorname{End}_{A}(M)\) 交换；因此，对每个 \(y \in M\)，我们有

\[
\begin{array}{l} u (y) = u \Big (\sum_ {i = 1} ^ {n} \langle x _ {i}, x _ {i} ^ {*} \rangle y \Big) = \sum_ {i = 1} ^ {n} u (\theta (x _ {i} ^ {*} \otimes y) (x _ {i})) \\ = \sum_ {i = 1} ^ {n} \theta (x _ {i} ^ {*} \otimes y) (u (x _ {i})) = \Big (\sum_ {i = 1} ^ {n} \langle u (x _ {i}), x _ {i} ^ {*} \rangle \Big) y. \end{array}
\]

因此，u 属于 \(A_{M}\)，从而 M 是平衡的。

推论 1. —— 自由模是平衡的。

当模为零时这是明显的，而非零自由模是生成模（VIII, 第 81 页，例 2）。

推论 2. —— 设 A 是一个环，n 是一个整数 \(\geqslant 0\)。\(\mathbf{M}_{n}(\mathrm{A})\) 的中心由元素属于 A 的中心的纯量矩阵组成。我们把 \(A^{n}\) 视为左 \(\mathbf{M}_{n}(\mathrm{A})\)-模（II, §10, No.~7, 第 349 页）。这个模的自同态是映射 \(x \mapsto xa\)，其中 a 取遍 A。

设 M 是右 A-模 \(A_{d}^{n}\)。由推论 1，它是平衡的。于是 \(A_{M}\)、\(A_{M}^{\prime}\) 与 \(A_{M}^{\prime\prime}\) 的中心重合。再由 \(A_{M}\) 与 A 等同且 \(A_{M}^{\prime}\) 与 \(\mathbf{M}_{n}(\mathbf{A})\) 等同，即得推论 2。

注. —— 设 M 是一个 A-模。若 \(A_{M}\)-模 M 是生成模，则我们有 \(A_{M} = A_{M}^{\prime\prime}\)（这由把定理 2 应用于 \(A_{M}\)-模 M 得到），从而 M 是平衡 A-模。

推论 3. —— 交换环上的每个有限生成投射模都是平衡的。

事实上，由 VIII, 第 82 页的命题 3，有限生成投射模 M 是生成 \(A_{M}\)-模。因此本推论由上面的注得出。

推论 4. —— 主理想整环上的每个有限生成模都是平衡的。

事实上，主理想整环 A 上的有限生成模 M 是生成 \(A_{M}\)-模（VIII, 第 81 页，例 4）。

推论 5. —— 设 K 是一个交换域（域），V 是 K 上一个有限维向量空间，u 和 v 是 V 的自同态。下列性质等价：

\begin{enumerate}
\def\labelenumi{(\roman{enumi})}
\item
  存在一个多项式 \(\mathrm{P}\)，它属于 \(\mathrm{K}[\mathrm{X}]\)，使得 \(v = \mathrm{P}(u)\)。
\item
  自同态 v 与 V 的每个与 u 交换的自同态交换。
\end{enumerate}

取 A 为环 K{[}X{]}，M 为由 V 和 u 得到的 K{[}X{]}-模（VII, §5, No.~1, 第 28 页）。断言 (i) 表示 \(v \in K[X]_{M}\)，(ii) 表示 \(v \in K[X]_{M}'\)。因此推论 5 是推论 4 的一个特例。

命题 4. —— 反模为有限生成的半单模是平衡的。

这由 VIII, 第 81 页的例 3 及该注得出。

推论 1. —— 设 \((\mathrm{S}_{i})_{i\in\mathrm{I}}\) 是一个由两两不同构的单 A-模组成的有限族。对 \(i\in I\)，用 \(D_{i}\) 表示 \(S_{i}\) 的自同态体之反环。假设对每个 \(i\in I\)，\(D_{i}\)-向量空间 \(S_{i}\) 是有限维的。则映射 \(a \mapsto (a_{\mathrm{S}_i})_{i \in \mathrm{I}}\)（从 A 到 \(\prod_{i \in \mathrm{I}} \mathrm{End}_{\mathrm{D}_i}(\mathrm{S}_i)\)）是满射。

考虑 A-模 \(M = \prod_{i \in I} S_i\)。由于 I 有限，我们又有 \(M = \bigoplus_{i \in I} S_i\)，且 \(S_i\) 在 M 中的像是 M 的 \(S_i\) 型同型分量。因此 A-模 M 的自同态是映射 \((s_i) \mapsto (s_i d_i)\)，其中 \((d_i)_{i \in I}\) 取遍 \(\prod_{i \in I} D_i\)（VIII, 第 66 页，命题 5）。由于 I 有限且对每个 \(i \in I\)，右 \(D_i\)-向量空间 \(S_i\) 是有限维的，M 的反模是有限生成的。由命题 4，A-模 M 是平衡的。而 A-模 M 的双交换子由 \(\text{End}_{\mathbf{Z}}(\text{M})\) 中形如 \(\prod_{i \in I} u_i\) 的元素组成，其中 \((u_i) \in \prod_{i \in I} \text{End}_{\text{D}_i}(\text{S}_i)\)（VIII, 第 78 页，命题 1），因为 \(\text{End}_{\text{D}_i}(\text{S}_i)\) 是 \(S_i\) 的双交换子。本推论由此得出。

特别地，当 A 是交换域（域）K 上的代数且每个 \(S_{i}\) 是作为 K-向量空间有限维的单 A-模时，本推论适用：事实上，\(D_{i}\) 这时含诸位似 \(\alpha_{S_{i}}\)，其中 \(\alpha\) 取遍 K，且 \(S_{i}\) 在 \(D_{i}\) 上有限维，因为它在 K 上如此。

推论 2（伯恩赛德定理）. —— 设 A 是代数闭交换域（域）K 上的代数，S 是一个作为 K-向量空间有限维的单 A-模。则映射 \(a \mapsto a_{\mathrm{S}}\)（从 A 到 \(\operatorname{End}_{\mathrm{K}}(\mathrm{S})\)）是满射。

事实上，A-模 S 的自同态体由诸位似 \(\alpha_{S}\)（\(\alpha \in K\)）组成（VIII, 第 47 页，定理 1）。于是我们对单 A-模 S 应用推论 1。

